%! Parametrization of Implicitly Defined Functions — Unit 4, Lesson 4.7 — Guided Notes
\documentclass[10pt]{article}
\usepackage{precalculus-article}
\usepackage{precalculus-boxes}
\newcommand{\TallMath}[1]{$\displaystyle #1\rule[-1.4em]{0pt}{3.2em}$}

\begin{document}

\pageheader{Unit 4, Lesson 4.7}{Guided Notes}
\namedateperiod

\vspace{0.1in}

\begin{objectivebox}
By the end of this lesson, I will be able to:
\begin{itemize}[leftmargin=1.5em, itemsep=3pt]
  \item \writeline
  \item \writeline
  \item \writeline
\end{itemize}
\end{objectivebox}

\vspace{0.08in}

\begin{vocabbox}
\termblanklong{implicit curve}
\termblanklong{parametrization}
\termblanklong{ellipse parametrization}
\termblanklong{parabola parametrization}
\end{vocabbox}

\vspace{0.08in}

\begin{hookbox}
The ellipse $\dfrac{x^2}{9}+\dfrac{y^2}{4}=1$ is \emph{not} a function --- it fails the vertical
line test. But if a particle traces this ellipse over time, we can use parametric tools to find
its position, extremes, and intercepts. How do we convert this implicit equation to a parametric form?

\writeline
\end{hookbox}

\vspace{0.08in}

\begin{notesbox}{Section 1 --- Ellipse Parametrization}
\small
\textbf{Setup.} Starting from $\dfrac{(x-h)^2}{a^2}+\dfrac{(y-k)^2}{b^2}=1$, set
$\dfrac{x-h}{a}=\cos t$ and $\dfrac{y-k}{b}=\sin t$.

Then: \quad $x(t) = $ \blank{4cm} \quad and \quad $y(t) = $ \blank{4cm}

Domain: $0 \le t \le $ \blank{1.5cm}

\vspace{0.08in}
\textbf{Verify.} Substitute back into the ellipse equation and use the identity
$\cos^2 t + \sin^2 t = 1$:

\vspace{0.06in}
\hspace{1em} $\dfrac{(x(t)-h)^2}{a^2}+\dfrac{(y(t)-k)^2}{b^2}
= \dfrac{(\blank{1.8cm})^2}{a^2}+\dfrac{(\blank{1.8cm})^2}{b^2}
= \blank{1cm} + \blank{1cm} = \blank{0.8cm}$ \quad \checkmark

\vspace{0.12in}
\textbf{Example.} Ellipse $\dfrac{x^2}{16}+\dfrac{y^2}{4}=1$.
Here $h=\blank{0.7cm}$, $k=\blank{0.7cm}$, $a=\blank{0.7cm}$, $b=\blank{0.7cm}$.

\vspace{0.06in}
Parametrization: \quad $x(t) = $ \blank{3.5cm} \quad $y(t) = $ \blank{3.5cm} \quad
$0 \le t \le \blank{1.2cm}$

\vspace{0.1in}
\textbf{Extrema from the parametrization.}
\renewcommand{\arraystretch}{1.6}
\begin{tabular}{lll}
Maximum $x(t)$: \blank{1.5cm} when $t = $ \blank{1.2cm} & $\Rightarrow$ & Point: \blank{2.2cm} \\
Minimum $x(t)$: \blank{1.5cm} when $t = $ \blank{1.2cm} & $\Rightarrow$ & Point: \blank{2.2cm} \\
Maximum $y(t)$: \blank{1.5cm} when $t = $ \blank{1.2cm} & $\Rightarrow$ & Point: \blank{2.2cm} \\
\end{tabular}
\end{notesbox}

\vspace{0.08in}

\begin{notesbox}{Section 2 --- Parabola Parametrization}
\small
\textbf{Setup.} For $y - k = a(x-h)^2$ (vertical parabola): set $x(t) = $ \blank{2.5cm},
then $y(t) = $ \blank{4cm}, $t \in \mathbb{R}$.

Vertex at $t=0$: \blank{2cm}

\vspace{0.1in}
\textbf{Example.} $y = x^2 - 4x + 3$.

Step 1 --- Complete the square: $y = (x - \blank{0.7cm})^2 - \blank{0.7cm}$.
\; Vertex: $(\blank{1cm},\,\blank{1cm})$.

Step 2 --- Parametrize: $x(t) = \blank{2cm}$, \quad $y(t) = \blank{3.5cm}$.

Step 3 --- Fill the table:

\renewcommand{\arraystretch}{1.6}
\begin{tabular}{|c|c|c|c|}
\hline
$t$ & $x(t)=2+t$ & $y(t)=t^2-1$ & Point \\
\hline
$-2$ & \blank{1.4cm} & \blank{1.4cm} & \blank{2cm} \\
\hline
$-1$ & \blank{1.4cm} & \blank{1.4cm} & \blank{2cm} \\
\hline
$\phantom{-}0$ & \blank{1.4cm} & \blank{1.4cm} & \blank{2cm} \\
\hline
$\phantom{-}1$ & \blank{1.4cm} & \blank{1.4cm} & \blank{2cm} \\
\hline
$\phantom{-}2$ & \blank{1.4cm} & \blank{1.4cm} & \blank{2cm} \\
\hline
\end{tabular}
\end{notesbox}

\vspace{0.08in}

\begin{notesbox}{Section 3 --- Applying Parametric Tools to Implicit Curves}
\small
\textbf{Ellipse:} $x(t) = 3\cos t$, \; $y(t) = 5\sin t$.

\begin{itemize}[leftmargin=1.4em, itemsep=4pt]
  \item When does the curve cross the \textbf{$y$-axis}?
    Set $x(t)=0$: $\cos t = 0$ $\Rightarrow$ $t = $ \blank{2.8cm}\\[2pt]
    \hspace{1.4em} Coordinates at those $t$-values: \blank{4cm}

  \item When does the curve cross the \textbf{$x$-axis}?
    Set $y(t)=0$: $\sin t = 0$ $\Rightarrow$ $t = $ \blank{2.8cm}\\[2pt]
    \hspace{1.4em} Coordinates at those $t$-values: \blank{4cm}

  \item What is the \textbf{rightmost} point? Max of $x(t)=3\cos t$ is \blank{0.8cm},
    occurring at $t = $ \blank{1cm}. Point: \blank{2cm}

  \item What is the \textbf{topmost} point? Max of $y(t)=5\sin t$ is \blank{0.8cm},
    occurring at $t = $ \blank{1cm}. Point: \blank{2cm}
\end{itemize}
\end{notesbox}

\end{document}
